% SPDX-License-Identifier: CC-BY-4.0
% Copyright (c) 2026 Martin Schröder <info@swedishembedded.com>
\section{Method}\label{sec:method}

Each generation repeats one cycle. The current champion proposes programs; the interpreter runs them and records their output.
The champion then answers each program several times with thinking enabled, and the interpreter checks every final answer.
A program enters the training set only where the champion's answers disagree with each other in correctness: sometimes right,
sometimes wrong. Problems the model always solves teach little, and problems it never solves give no verified trace to learn
from; the same rule is what makes a group of sampled answers informative for policy-gradient training (dynamic sampling
in~\cite{yu2025dapo}). For each admitted program the shortest correct reasoning trace is kept, in the spirit of self-taught
reasoning~\cite{zelikman2022star} and self-training on verified samples~\cite{singh2023rest}. A fresh adapter is trained from
the base on all verified traces kept so far, and it replaces the champion only if it passes a gate on the learner's own
development tasks; the gate may reject, and a rejected candidate is recorded with its evidence. The learner hands each new
champion, with the list of programs it trained on, to the examiner, and never receives anything back except signed aggregate
verdicts.
